{
\clearpage
\phantomsection
\addcontentsline{toc}{sectionnoiddot}{《本科生毕业论文（设计）任务书》}
}

\begin{center}
  \setfont{-2}{1.5}\bfseries 本科生毕业论文（设计）任务书
\end{center}

\blankline{5}{1}

\noindent {\setfont{-4}{1.5}\bfseries 一、题目：\ProjectName}\par
\noindent {\setfont{-4}{1.5}\bfseries 二、指导教师对毕业论文（设计）的进度安排及任务要求}\par
%----------正文----------

\noindent 任务要求：

（1）阅读论文"Y. Zhang, et al, Hand Gesture Recognition for Smart Devices by Classifying Deterministic Doppler Signals, \textit{IEEE Transactions on Microwave Theory and Techniques}, vol. 69, no. 1, pp. 365--377, Jan. 2021"。

（2）研读相关文献，掌握多普勒效应的原理及其在连续波雷达中的具体应用。

（3）使用实验室的多普勒雷达传感器，完成对两种及以上不同手势的测量。

（4）构建机器学习模型对测得的手势数据进行分类和识别，评估模型的性能。

（5）完成毕业论文（设计）的总结工作和毕业论文的撰写（11000字以上）。

\vspace{1ex}
\noindent 进度安排：

\noindent 2025-11-10至2025-11-16：明确毕设任务、目标和进度要求。

\noindent 2025-11-17至2025-11-30：查阅文献资料，调研国内外研究现状，撰写文献综述文档；完成外文翻译，撰写文献翻译文档。

\noindent 2025-12-01至2025-12-14：学习实验室多普勒雷达传感器的使用，学会基础的数据采集和信号处理。

\noindent 2025-12-15至2025-12-21：撰写开题报告，应包含学生本人已开展的阶段性工作及结果。

\noindent 2025-12-22至2025-12-29：讨论修改文献综述、外文翻译和开题报告，准备开题答辩。

\noindent 2026-01-05至2026-01-11：开题答辩。

\noindent 2026-03-02至2026-03-22：采用连续波多普勒传感器（工作频率5.8\,GHz），完成手势运动的采集和测量。

\noindent 2026-03-23至2026-03-27：中期检查。

\noindent 2026-03-28至2026-04-12：设计相应的机器学习模型，对采集到的手势数据进行分类识别。

\noindent 2026-04-13至2026-04-20：撰写、上传毕业论文。

\noindent 2026-04-21至2026-05-24：答辩前审核流程。

\noindent 2026-05-25至2026-05-29：论文答辩。

%----------正文----------
\vfill
{\setfont{-4}{1.5}\bfseries
\noindent 起讫日期\hspace{1em}2025年11月10日至2026年5月29日\vspace{0.5ex}\par
\hfill 指导教师（签名）{\underline{\makebox[7.5em][c]{}}}\hspace{0.5em}职称{\underline{\makebox[6.5em][c]{教授}}}\vspace{0.5ex}
}

\noindent {\setfont{-4}{1.5}\bfseries 三、系或研究所审核意见}\par
\multido{}{3}{\blankline{-4}{1.5}}
{\setfont{-4}{1.5}\bfseries
\hfill 负责人（签名）{\underline{\makebox[8em][c]{}}}\par
\hfill 2026年XX月XX日
}